
    \section{Convex Hull}
        \subsection{Rotation Calipers}
        Diameter from a convex Polygon.
        
        Solution $O(n)$. Original Problem (NCPC 2013 - D).
        
        \ \\
        
        \begin{tabular}{|p{31cm}|}
          \begin{verbatim}
            Implementacao do Convex Hull do CP-ALGORITHMS
            Para usar crie um vetor com os pontos,
            		vector<pt> pts;
            e execute
            		convex_hull(|pts|);
            Apos isso, o vector pts sera alterado e ele so tera os pontos do Convex Hull 
            na ordem horaria, comecando do elemento de menor x, como segunda condicao menor y.
            
            Como funciona o algoritmo
            - Acha os dois extremos de x.
            - Monta dois subgrupos, os "up" e os "down", em relacao
            Convex Hull na parte superior e inferior.
            - Itera pelos pontos e usa o produto vetorial para ver se
            ele forma um sentido horario, se formar adiciona.
            - Remove todos os anteriores ao ultimo ponto adicionado que agora, com o ponto
            atual, formando um sentido anti-horario
            \end{verbatim}
        \end{tabular}
        \\
        
        \inputwithhash{./solutions/Geometry/MaxDistanceBetweenPoints.cpp}
